\section{Binary-Outcome Calibration: Design and Results}\label{app:binary-calibration}
Clark's higher-education outcome is a binary indicator observed for men in
these historical cohorts (the supplementary methods, p.~14, describe the
male-only construction; the two cohorts and their pair totals appear in
Table~2, p.~2, of \citet{Clark2023}, and in Table~4.5, p.~54, of
\citet{Clark2026}). We recovered cohort prevalences from the archived
supplementary workbook `pnas.2300926120.sd03.xlsx`, sheet `Table 2 Ded
1780-1919`, rather than from a published marginal table. After deduplicating
person IDs and retaining nonmissing male indicators, the fractions coded one
were $432/16{,}238=0.02660$ for births in 1780--1859 and
$522/21{,}150=0.02468$ for births in 1860--1919. These are workbook-based
calibrations, not claims that every kinship-pair sample has exactly those
marginals.

The simulation is fully specified as follows. For each of the 11 kinship rows
in the Fisher table (the $n,D$ structure is Table~A1, p.~296, of
\citet{Clark2026}), we generated $1{,}000{,}000$ independent pairs of latent
standard-normal phenotypes with correlation
\[
 \rho_{n,D}=h^2\left(\frac{1+m}{2}\right)^n
 \left(\frac{1+r}{1+m}\right)^D.
\]
The generating values $h^2=0.60$, $r=0.80$, and $m=h^2r=0.48$ are deliberately
chosen illustrative values, not estimates read from Clark. For each endpoint
we then formed the standardized recorded continuous outcome
$Y=\sqrt{\theta}P+\sqrt{1-\theta}U$, with independent standard-normal $U$ and
$\theta=0.80$; thus 20\% of recorded variance is independent white noise.
For cohort prevalence $p$, we defined
$\mathrm{Edu}=1\{Y\geq\Phi^{-1}(1-p)\}$, so the threshold was applied to the
noisy observed outcome and the population prevalence was exactly $p$. The
random seed was 20260920; pairs were simulated separately by kinship, not as
one connected pedigree.

For each simulated dataset we minimized the equally weighted sum of squared
residuals in the 11 log correlations, imposing the phenotypic-assortment
restriction $m=h^2r$. The fixed-$\theta$ fit set $\theta=1$; the free-$\theta$
fit estimated $(h^2,r,\theta)$ in $(0,1]^3$. The table reports the resulting
estimates; the displayed fixed-$\theta$ rows use the admissible bounds
$h^2,r\leq1$.

\begin{center}
\small
\begin{tabular}{llrrrrr}
\hline
Cohort & Outcome & $\theta$ & $h^2$ & $r$ & $m$ & RMSE(log)\\
\hline
1780--1859 & Continuous & 1 fixed & .476 & 1.000 & .476 & .044\\
1780--1859 & Continuous & free & .606 & .782 & .474 & .009\\
1780--1859 & Binary Edu & 1 fixed & .227 & 1.000 & .227 & .363\\
1780--1859 & Binary Edu & free & .412 & .906 & .373 & .072\\
1860--1919 & Continuous & 1 fixed & .480 & 1.000 & .480 & .048\\
1860--1919 & Continuous & free & .598 & .809 & .484 & .012\\
1860--1919 & Binary Edu & 1 fixed & .226 & 1.000 & .226 & .394\\
1860--1919 & Binary Edu & free & .394 & .996 & .392 & .071\\
\hline
\end{tabular}
\end{center}

The continuous free-$\theta$ fits recover the deliberately chosen generating
values closely. The binary free-$\theta$ fits do not: at Clark-like education
prevalences they understate $h^2$ by about 31--34\%, inflate $r$ from .80 to
.91--1.00, estimate $m$ around .37--.39 instead of .48, and estimate
$\theta$ around .30 instead of .80. Fixing $\theta=1$ is worse and reaches
$r=1$. Thresholding a rare binary outcome therefore changes the correlation
pattern across kinships, rather than applying a common multiplicative
attenuation. These are calibrated simulations, not a formal rejection of
Clark's data; they show why the reported binary Pearson correlations need an
explicit threshold/liability model before the continuous Fisher formulas can
be given a latent interpretation.
